\documentclass{article}
\usepackage{booktabs}
\usepackage[a4paper]{geometry}
\usepackage{fancyhdr}
\pagestyle{fancy}
\lhead{Soziale Ungleichheit (in Deutschland)}
\rhead{März 2026}
\begin{document}
\section{Soziale Ungleichheit (in Deutschland)}
In der Realität ist die soziale Marktwirtschaft nur in geringen Teilen in der Lage die soziale Ungleichheit der Marktwirtschaft zu bekämpfen. Diese ist in Deutschland nach wie vor sehr groß.
 
Daraus folgt eine fehlende Chancengleichheit. 
 
\subsection{Gini-Koeffizient} 
Mithilfe einer Lorenzkurve kann die (un)gleichheit der Verteilung von Einkommen oder Vermögen dargestellt werden. Auf der $x$-Achse liegt dabei der Prozentanteil der Bevölkerung ($0\% - 100\%$), auf der $y$-Achse liegt prozentuale Anteil des Vermögens den die ärmste Anteil der Bevölkerung der $x$-Achse hat.
  
Hätten alle gleichviel, so läge die Linie in einem $45^\circ$ Winkel.
 
Die tatsächliche Linie liegt aber immer unter dieser hypothetischen gerechtesten Verteilung; also wird das Integral von dessen Differenz genommen und durch das Integral der gerechtesten Verteilung geteilt. Dies ist der \emph{Gini}-Koeffizient.
 
Ist er 0, so gibt es keine Differenz zwischen der gerechtesten Verteilung und der echten Verteilung, ist er 1, so gibt es keine Differenz zwischen der ungerechtesten Verteilung und der echten Verteilung.
 
Siehe unten für die zurzeitigen Werte in Deutschland.
 
\subsection{Weitere Indikatoren}
Desweiteren können andere Statistiken angemerkt werden, wie dass die reichsten $10\%$ ungefähr $60\%$ des Vermögens besitzen oder dass die Kinder von NichtakademikerInnen selbst seltener AkademikerInnen werden.
 
\subsection{Entwicklung}
Seit 2010 sind innerhalb von zehn Jahren sowohl die Einkommens- als auch die Vermögensverteilung ungerechter geworden.
\begin{center}
 \begin{tabular}{lll}
  \multicolumn{3}{c}{Gini-Koeffizienten} \\
  \toprule 
  & 2010 & 2020 \\
  \midrule
  Einkommen & 0,293 & 0,309 \\
  Vermögen  & 0,684 & 0,816 \\
  \bottomrule
 \end{tabular}
\end{center}
 
\end{document}